% =====================================================
% CAPA / ABERTURA
% =====================================================
\capaTutorial{\faGlobe}{Registro de Domínio Próprio Gratuito \& Deploy Web}
  {Guia prático para a Estação Meteorológica e projetos web/IoT com GitHub Student Pack}
  {\chip{\faCloud}{Estação Meteorológica}\hspace{4pt}
   \chip{\faGithub}{GitHub Pages \& DNS}\hspace{4pt}
   \chip{\faGraduationCap}{GitHub Student Pack}}

\vspace{0.4em}
\cartoesUC{Domínio próprio oficial registrado sem custos e conectado com HTTPS ao dashboard da equipe}

\begin{boxObjetivos}
\begin{itemize}
  \item \textbf{Ativar o benefício do GitHub Student Developer Pack} utilizando o e-mail institucional do IFSC (\texttt{@aluno.ifsc.edu.br});
  \item \textbf{Resgatar um domínio de primeiro nível gratuito por 1 ano} (\texttt{.me} ou \texttt{.tech}) sem necessidade de cartão de crédito;
  \item \textbf{Configurar os apontamentos de DNS (CNAME e registros A)} direcionados aos servidores do GitHub Pages;
  \item \textbf{Ativar o domínio personalizado com HTTPS forçado} no repositório do projeto, garantindo 99,9\% de disponibilidade para bancas e avaliadores.
\end{itemize}
\end{boxObjetivos}

\section{Contexto \& Arquitetura Recomendada}

Para sistemas que integram hardware e web (como a Estação Meteorológica com Raspberry Pi e sensor BME280), manter o servidor web dependente exclusivamente da conexão local da placa pode trazer instabilidade (quedas de energia, IP dinâmico ou bloqueios de firewall/CGNAT da operadora).

\begin{boxAtencao}[Por Que Evitar No-IP e Antigos Domínios Gratuitos?]
Historicamente, muitos estudantes utilizam serviços como o \textbf{No-IP} gratuito ou domínios do antigo \textbf{Freenom} (\texttt{.tk}, \texttt{.ml}). Ambos devem ser evitados:
\begin{itemize}
  \item O \textbf{Freenom faliu e foi descontinuado} internacionalmente em decorrência de litígios jurídicos e abuso de spam;
  \item O plano gratuito do \textbf{No-IP exige renovação manual por clique em e-mail a cada 30 dias} com captcha. Se a equipe esquecer, o link sai do ar repentinamente.
\end{itemize}
A solução profissional adotada pela indústria acadêmica é resgatar um \textbf{domínio próprio oficial de primeiro nível} através do programa educacional do GitHub.
\end{boxAtencao}

\begin{destaque}{ciano}{\faLayerGroup}{Arquitetura em Duas Camadas (Edge + Cloud)}
A abordagem mais moderna e resiliente para o projeto é o desacoplamento:
\begin{enumerate}
  \item \textbf{Camada de Coleta (Raspberry Pi):} Script em Python realiza as leituras periódicas do sensor BME280 e apenas envia os dados (via requisição POST simples) para um banco/API em nuvem;
  \item \textbf{Camada de Visualização (Dashboard Web):} O painel em HTML5/JavaScript fica hospedado na infraestrutura de alta disponibilidade do \textbf{GitHub Pages}, acessível publicamente através do domínio personalizado com certificado SSL automático.
\end{enumerate}
\end{destaque}

\section{Passo a Passo: Da Ativação ao Ar}

\passo{1}{Vincular o E-mail Institucional e Ativar o GitHub Student Pack}

Qualquer uma das integrantes da equipe matriculadas no curso técnico integrado pode realizar este procedimento:

\begin{enumerate}
  \item Acesse sua conta no GitHub e navegue até \textbf{Settings $\rightarrow$ Emails};
  \item No campo \textit{Add email address}, adicione seu e-mail institucional do IFSC (ex.: \texttt{nome.sobrenome@aluno.ifsc.edu.br}) e confirme a mensagem de verificação recebida na sua caixa de entrada;
  \item Em seguida, abra o portal do programa: \url{https://education.github.com/pack};
  \item Clique no botão verde \textbf{Sign up for Student Developer Pack} e selecione o seu e-mail \texttt{@aluno.ifsc.edu.br};
  \item Como o domínio do IFSC já é pré-aprovado na base acadêmica do GitHub Education, a validação é concluída rapidamente (muitas vezes de forma automática).
\end{enumerate}

\passo{2}{Resgatar o Domínio Gratuito (Namecheap ou .TECH)}

Com o pacote ativado, a equipe tem direito a um domínio de primeiro nível gratuito por 12 meses:

\begin{destaque}{verdeVivo}{\faGlobe}{Opção A: Domínio .ME via Namecheap (nc.me) — Recomendada}
\begin{enumerate}
  \item Acesse diretamente o portal estudantil da Namecheap: \url{https://nc.me};
  \item Clique no botão superior direito \textbf{Sign in with GitHub} e autorize a integração;
  \item No campo de busca, pesquise o domínio desejado para o projeto (veja as sugestões abaixo);
  \item O valor do carrinho será atualizado automaticamente para \textbf{\$0.00} (100\% gratuito com WhoisGuard e SSL);
  \item Conclua o cadastro na Namecheap para ter acesso ao painel de gerenciamento de DNS.
\end{enumerate}
\end{destaque}

\begin{boxDica}[Nomes de Domínio Testados e Disponíveis para Registro]
Testamos a disponibilidade DNS no momento da elaboração deste material e os seguintes endereços estão livres para registro imediato:
\begin{itemize}
  \item \texttt{estacaogaropaba.me}
  \item \texttt{estacao-ifsc.me}
  \item \texttt{estacaometeorologica-ifsc.me}
  \item \texttt{estacaogaropaba.tech} (disponível via \href{https://get.tech/github-student-pack}{\texttt{get.tech/github-student-pack}})
\end{itemize}
\end{boxDica}

\passo{3}{Configurar os Apontamentos de DNS para o GitHub Pages}

Para que os navegadores saibam que o domínio registrado deve abrir o site hospedado no GitHub, é necessário criar registros de DNS no painel da registradora (Namecheap ou .TECH):

\begin{enumerate}
  \item Entre no painel da Namecheap, acesse \textbf{Domain List $\rightarrow$ Manage $\rightarrow$ Advanced DNS};
  \item Remova os registros padrão existentes e adicione os seguintes apontamentos:
\end{enumerate}

\begin{table}[H]
\centering
\small
\begin{tabularx}{\textwidth}{l l X l}
\toprule
\textbf{Tipo de Registro} & \textbf{Host} & \textbf{Valor / Destino} & \textbf{TTL} \\
\midrule
\texttt{CNAME} & \texttt{www} & \texttt{sophirie.github.io} & Automático \\
\texttt{A} & \texttt{@} & \texttt{185.199.108.153} & Automático \\
\texttt{A} & \texttt{@} & \texttt{185.199.109.153} & Automático \\
\texttt{A} & \texttt{@} & \texttt{185.199.110.153} & Automático \\
\texttt{A} & \texttt{@} & \texttt{185.199.111.153} & Automático \\
\bottomrule
\end{tabularx}
\caption{Registros DNS oficiais dos servidores do GitHub Pages.}
\end{table}

\passo{4}{Ativar o Domínio Personalizado no Repositório do GitHub}

Agora basta informar ao repositório do projeto qual é o domínio que ele deve atender:

\begin{enumerate}
  \item Acesse o repositório da equipe: \url{https://github.com/sophirie/Estacao-Meteorologica-no-IFSC-Garopaba};
  \item Clique na aba \textbf{Settings $\rightarrow$ Pages};
  \item No campo \textbf{Custom domain}, digite o domínio registrado (ex.: \texttt{estacaogaropaba.me}) e clique em \textbf{Save};
  \item O GitHub criará automaticamente um arquivo chamado \texttt{CNAME} na raiz do repositório;
  \item Aguarde alguns instantes pela checagem de DNS e marque a caixa \textbf{Enforce HTTPS}. O GitHub emitirá um certificado SSL gratuito (Let's Encrypt) automaticamente.
\end{enumerate}

\section{Validação \& Verificação Técnica}

\begin{boxTeste}[Como Testar a Propagação do Domínio]
Após salvar as configurações de DNS, o processo de propagação mundial leva geralmente de 5 a 30 minutos. Para conferir o status:
\begin{enumerate}
  \item \textbf{Pelo Terminal:}
\begin{saida}[Testando a resolução DNS no terminal]
# Consulta o registro CNAME do subdomínio www
nslookup www.estacaogaropaba.me

# Consulta os registros A do domínio raiz
dig +short estacaogaropaba.me
# Resposta esperada: os IPs 185.199.108.153 ... do GitHub
\end{saida}
  \item \textbf{Pelo Navegador:} Abra uma janela anônima e acesse \url{https://estacaogaropaba.me}. O cadeado de conexão segura deve aparecer verde/fechado ao lado da URL.
\end{enumerate}
\end{boxTeste}

\begin{boxChecklist}[Checklist de Entrega para a Equipe]
\begin{itemize}[label={\faSquare[regular]}]
  \item Solicitação aprovada no GitHub Student Developer Pack com e-mail do IFSC;
  \item Domínio registrado na Namecheap (\texttt{nc.me}) ou .TECH com custo \$0.00;
  \item Registros de DNS (\texttt{CNAME} e 4 registros \texttt{A}) configurados;
  \item Opção \textit{Enforce HTTPS} ativa no GitHub Pages;
  \item URL final atualizada no card do projeto e no artigo SBC no Overleaf.
\end{itemize}
\end{boxChecklist}
